\documentclass[11pt]{article}
\usepackage[margin=1in]{geometry}
\usepackage{amsmath,amssymb,amsthm,mathtools}
\usepackage{booktabs,array,enumitem,microtype}
\usepackage[hidelinks]{hyperref}
\hypersetup{
  pdftitle={The Complexity of Stationary Diffusion-Driven Instability in Mass-Action Reaction Networks},
  pdfauthor={Alec Kriebel},
  pdfsubject={Exact semialgebraic characterization, NP-hardness, and fixed-species algorithms},
  pdfkeywords={mass-action kinetics, Turing instability, reaction networks, NP-hardness, real algebraic geometry}
}

\usepackage[nameinlink,capitalise]{cleveref}
\usepackage{xcolor}

\newtheorem{theorem}{Theorem}[section]
\newtheorem{proposition}[theorem]{Proposition}
\newtheorem{lemma}[theorem]{Lemma}
\newtheorem{corollary}[theorem]{Corollary}
\theoremstyle{definition}
\newtheorem{definition}[theorem]{Definition}
\newtheorem{remark}[theorem]{Remark}
\newcommand{\R}{\mathbb R}
\newcommand{\Q}{\mathbb Q}
\newcommand{\diag}{\operatorname{diag}}
\newcommand{\im}{\operatorname{im}}
\newcommand{\ri}{\operatorname{ri}}
\newcommand{\rank}{\operatorname{rank}}
\newcommand{\Hur}{\operatorname{Hur}}

\title{The Complexity of Stationary Diffusion-Driven Instability\\
\large Exact Semialgebraic Characterization, NP-Hardness, and Fixed-Species Algorithms}
\author{Alec Kriebel\\\small \texttt{me@aleckriebel.com}}
\date{August 2026}

\begin{document}
\maketitle

\begin{abstract}
We consider a finite classical mass-action reaction network specified by its
indexed source-complex matrix $Y$ and stoichiometric matrix $\Gamma$.  We ask
whether some strictly positive rates and homogeneous equilibrium are stable
relative to the stoichiometric compatibility class, while some strictly
positive all-species diagonal diffusion matrix creates a nonzero spatial mode
with a positive real eigenvalue.  We prove an exact characterization: there
must exist a positive steady flux $v$, $\Gamma v=0$, and a positive reciprocal
concentration vector $h$ such that
$J=\Gamma\diag(v)Y^T\diag(h)$ is Hurwitz on $\im\Gamma$ and has a negative
signed principal minor.  This yields a polynomial-size formula in the
existential theory of the reals and finite exact certificates on both sides.
The decision problem is NP-hard under a polynomial-time reduction from
\textsc{Partition}.  For every fixed species count, however, it is decidable
in polynomial bit complexity by projecting the positive flux cone into the
fixed-dimensional space of Jacobian factors.  Conservation laws are handled
without restricting nonzero diffusive modes to the stoichiometric subspace.
The hardness construction uses reaction complexes of unbounded molecularity
and proves only weak NP-hardness.  The results concern weak stationary linear
instability; they do not assert a first or transverse crossing, exclude an
earlier wave instability, or prove a nonlinear patterned state.
\end{abstract}

\section{Introduction}

Turing's diffusion-driven instability mechanism asks how diagonal diffusion
can destabilize a homogeneous equilibrium that is stable without spatial
transport \cite{Turing1952}.  Matrix criteria are well developed when the
Jacobian is fixed \cite{Satnoianu2000,WangLi2001,Anma2012,Piskovsky2024}, and
interaction-topology studies have identified structural features and surveyed
small networks \cite{Marcon2016,Diego2018,Scholes2019}.  A reaction network,
however, does not specify one Jacobian.  Under classical mass action, positive
rates and positive equilibria generate a constrained family whose entries are
coupled by steady flux and concentration scaling.

This paper studies the corresponding network-wide decision problem.  The
input is the complete indexed pair $(\Gamma,Y)$, not merely a signed
interaction graph.  We call the network \emph{weakly all-mobile stationary
Turing-admissible} if there are strictly positive rate constants, a strictly
positive homogeneous equilibrium, and strictly positive species diffusion
coefficients such that the homogeneous Jacobian is Hurwitz on the
stoichiometric subspace and a nonzero Neumann mode has a strictly positive real
eigenvalue.

Our first contribution is an exact reduction of this question to positive
steady flux, relative Routh--Hurwitz inequalities, and one negative signed
principal minor.  The diffusion coefficients, domain scale, original rates,
and equilibrium concentrations are eliminated from the destabilization
condition.  Boolean selector variables encode the principal-set disjunction
with polynomial size, proving membership in the existential theory of the
reals.

The second contribution is a complexity boundary.  Starting from the
open-interval matrix-stability construction of Blondel and Tsitsiklis
\cite{BlondelTsitsiklis1997}, we introduce a row-splitting lift whose stability
is unchanged and whose arbitrary positive right scalings can be eliminated by
an inertia and determinant-continuity argument.  We then realize the entire
lifted row family---not just one matrix---as the positive steady-flux image of
a classical mass-action network.  This gives weak NP-hardness from
\textsc{Partition}.  The reduction has full stoichiometric rank and unit
stoichiometric columns, but the source and target molecularities are not
uniformly bounded.

The third contribution is tractability at fixed species count.  Extreme rays
of the nonnegative steady-flux cone are positive circuits supported on at most
$n+1$ reactions.  For fixed $n$ they can be enumerated in polynomial time,
projected into $\R^{n\times n}$, and converted in fixed dimension to a
relative-facet description.  Fixed-variable real-algebraic decision then gives
a polynomial-time exact algorithm for every fixed $n$, including arbitrary
one-, two-, and three-species reaction lists.

\paragraph{Scope.}
The word \emph{stationary} here means that the unstable spatial matrix has a
positive real eigenvalue.  We do not require this to be the first loss of
stability along a diffusion path, do not require simplicity or transversality,
and do not exclude an earlier complex crossing.  A linear Turing instability
also does not by itself guarantee a persistent nonlinear pattern
\cite{Krause2024}.  No such conclusion is made.

\paragraph{Relation to structural reaction-network work.}
Graph criteria for zero-eigenvalue Turing instability under mass action
\cite{MinchevaCraciun2013}, unstable cores and child selections under
parameter-rich kinetics \cite{VassenaStadler2024,Vassena2025}, and recent
sufficient conditions for networks with monomial steady-state
parametrizations \cite{Conradi2026} are closest in subject.  Their kinetic
classes, quantifiers, or necessity/sufficiency scopes differ from the
arbitrary-network classical-mass-action decision problem considered here.
The fixed-matrix principal-minor result is used in an exact form and re-proved;
it is not presented as the principal novelty.

\section{Mass-action Jacobians and conservation laws}
\label{sec:setup}

Consider indexed reactions
\[
 y_r\longrightarrow y'_r,\qquad r=1,\ldots,m,
\]
with $y_r,y'_r\in\mathbb Z_{\geq0}^n$ and $y_r\ne y'_r$.  Let
\[
 Y=[y_1\ \cdots\ y_m],\qquad
 \Gamma=[y'_1-y_1\ \cdots\ y'_m-y_m],\qquad
 S=\im\Gamma.
\]
With positive rates $k_r$, the vector field is
\[
 f_k(x)=\Gamma r_k(x),\qquad r_{k,r}(x)=k_rx^{y_r}.
\]

\begin{proposition}[Positive-equilibrium parametrization]
\label{prop:param}
At a positive equilibrium $x^*$ define
\[
 v_r=k_r(x^*)^{y_r},\qquad h_i=(x_i^*)^{-1}.
\]
Then $v>0$, $\Gamma v=0$, and
\begin{equation}
 J(v,h)=\Gamma\diag(v)Y^T\diag(h).
 \label{eq:Jfactor}
\end{equation}
Conversely, every $v>0$ with $\Gamma v=0$ and every $h>0$ are realized by
\[
 x_i^*=h_i^{-1},\qquad k_r=\frac{v_r}{(x^*)^{y_r}}.
\]
\end{proposition}

\begin{proof}
The derivative of $k_rx^{y_r}$ at $x^*$ is
$v_ry_r^T\diag(h)$, giving \eqref{eq:Jfactor}.  The equilibrium condition is
$\Gamma v=0$.  The displayed converse rates are positive, recover the flux
$v$, and therefore give both equilibrium and the same Jacobian.
\end{proof}

This parametrization includes zero complexes, parallel indexed reactions,
autocatalytic reactions, conservation laws, and arbitrary molecularity.

\begin{proposition}[Conservation splitting]
\label{prop:conservation}
Suppose $J(\R^n)\subseteq S$, $J(S)\subseteq S$, and $J|_S$ is Hurwitz.  Then
\[
 \R^n=S\oplus\ker J,
 \qquad
 J\sim\begin{pmatrix}J|_S&0\\0&0\end{pmatrix}.
\]
The conservation-induced zero eigenvalues are semisimple and, with
$s=\dim S$,
\begin{equation}
 \det(\lambda I-J)=\lambda^{n-s}\det(\lambda I_s-J|_S)>0
 \qquad(\lambda>0).
 \label{eq:positive-axis}
\end{equation}
\end{proposition}

\begin{proof}
The restriction is invertible, so $J(S)=S$.  Since the full image is contained
in $S$, $\im J=S$.  Invertibility on $S$ gives $S\cap\ker J=\{0\}$, and
rank--nullity gives the direct sum.  Both summands are invariant and $J$
vanishes on its kernel.  The determinant factorization and positivity follow.
\end{proof}

Choose a rational basis matrix $B\in\Q^{n\times s}$ for $S$ and a rational
left inverse $C$ with $CB=I_s$.  Then
\begin{equation}
 R(v,h)=CJ(v,h)B
 \label{eq:reduced}
\end{equation}
represents $J|_S$.  If $\widetilde B=BT$, the new representation is
$T^{-1}R(v,h)T$; all Hurwitz tests are basis-independent.

For Neumann boundary conditions, a nonconstant Laplacian eigenfunction
$\phi$ has zero spatial mean.  Therefore a perturbation $a\phi(x)$ preserves
every integrated conservation law for every amplitude $a\in\R^n$.  The
nonzero-mode amplitude space is the full species space, and its matrix is
$J-q^2D$.  It would generally be wrong to restrict this matrix to $S$: unequal
diagonal diffusion need not preserve $S$ pointwise.

\section{Fixed-J diagonal damping and the exact characterization}
\label{sec:characterization}

For $I\subseteq[n]$, write
\[
 c_I(J)=(-1)^{|I|}\det J[I,I],\qquad c_\varnothing(J)=1.
\]

\begin{theorem}[Positive diagonal damping]
\label{thm:fixedJ}
For every real square matrix $J$, the following are equivalent:
\begin{enumerate}[label=(\roman*)]
\item some $D\succ0$ diagonal makes $J-D$ have a strictly positive real
eigenvalue;
\item $c_I(J)<0$ for some principal set $I$.
\end{enumerate}
If $J$ is Hurwitz, a witnessing set is nonempty and proper.
\end{theorem}

\begin{proof}
Multilinearity in the columns gives
\begin{equation}
 \det(D-J)=\sum_{I\subseteq[n]}c_I(J)\prod_{j\notin I}d_j.
 \label{eq:principalexpansion}
\end{equation}
If every $c_I(J)\geq0$, then for $\lambda>0$ the same expansion with
$d_j+\lambda$ is strictly positive because the empty-set term is positive.
Thus $\det(\lambda I-(J-D))$ cannot vanish at a positive real $\lambda$.

Conversely, fix $I$ with $c_I(J)<0$ and set $d_j=t$ off $I$ and $d_j=t^{-1}$
on $I$.  The exponent of the $I$-term exceeds that of the $K$-term by
$|I\triangle K|$, so the negative $I$-coefficient uniquely dominates for
large finite $t$.  Hence $\det(D-J)<0$.  The monic characteristic polynomial
of $J-D$ is negative at zero and tends to $+\infty$ on the positive real
axis, so it has a positive real root.  When $J$ is Hurwitz,
$c_\varnothing=1$ and $c_{[n]}=\det(-J)>0$.
\end{proof}

Set
\[
 A(v)=\Gamma\diag(v)Y^T.
\]
Column scaling gives, for every principal set,
\begin{equation}
 (-1)^{|I|}\det J(v,h)[I,I]
 =\left(\prod_{i\in I}h_i\right)(-1)^{|I|}\det A(v)[I,I].
 \label{eq:minoridentity}
\end{equation}
Equivalently, Cauchy--Binet expands the second factor as
\[
 (-1)^{|I|}
 \sum_{\substack{T\subseteq[m]\\|T|=|I|}}
 \det(\Gamma_{I,T})\det(Y_{I,T})\prod_{r\in T}v_r.
\]
Cancellations among reaction selections are therefore retained exactly.

\begin{theorem}[Exact network characterization]
\label{thm:maincharacterization}
A finite classical mass-action network is weakly all-mobile stationary
Turing-admissible if and only if there exist
\[
 v\in\R_{>0}^m,
 \qquad h\in\R_{>0}^n,
 \qquad \varnothing\ne I\subseteq[n],
\]
such that
\begin{align}
 \Gamma v&=0,\label{eq:flux}\\
 R(v,h)&\text{ is Hurwitz},\label{eq:RH}\\
 (-1)^{|I|}\det A(v)[I,I]&<0.\label{eq:negative}
\end{align}
A witnessing set is necessarily proper.
\end{theorem}

\begin{proof}
Given a positive mass-action equilibrium, \cref{prop:param} produces $v,h$.
Homogeneous stability relative to the compatibility class is exactly
\eqref{eq:RH}.  If a nonzero mode has a positive real eigenvalue,
\cref{thm:fixedJ} gives a negative signed principal minor of $J$; its sign is
that of \eqref{eq:negative} by \eqref{eq:minoridentity}.

Conversely, reconstruct $x^*$ and $k$ from $v,h$.  The reduced matrix is
Hurwitz by \eqref{eq:RH}.  Equations \eqref{eq:minoridentity} and
\eqref{eq:negative} give a negative signed principal minor of $J$, so
\cref{thm:fixedJ} supplies $D\succ0$ for which $J-D$ has a positive real
eigenvalue.  On $(0,\pi)$ this is the first nonzero Neumann mode; other domain
scales are absorbed into $D$.  Conservation laws permit every nonzero-mode
amplitude, as shown in \cref{sec:setup}.  The full set cannot be negative:
its determinant is zero in the conservative case and has positive signed
determinant in full rank.
\end{proof}

Let the characteristic polynomial of $R(v,h)$ be
$\lambda^s+a_1\lambda^{s-1}+\cdots+a_s$.  The Routh--Hurwitz determinants
$\Delta_1,\ldots,\Delta_s$ turn \eqref{eq:RH} into strict polynomial
inequalities.

\begin{corollary}[Existential-real membership]
\label{cor:etrmembership}
The arbitrary-network decision problem belongs to the existential theory of
the reals and has a polynomial-size formula computed from $(\Gamma,Y)$.
\end{corollary}

\begin{proof}
Introduce Boolean variables $b_i^2-b_i=0$ and
$S_b=\diag(b_1,\ldots,b_n)$.  If $I=\{i:b_i=1\}$, simultaneous permutation of
rows and columns gives
\begin{equation}
 \det(I-S_b-S_bA(v)S_b)=(-1)^{|I|}\det A(v)[I,I].
 \label{eq:selector}
\end{equation}
The empty selector gives $1$ and cannot satisfy negativity.  Determinants,
characteristic coefficients, and Hurwitz determinants have polynomial-size
arithmetic circuits.  Rational row reduction produces $B,C$ with polynomial
encoding length, and gate variables convert the circuits to a polynomial-size
existential-real formula.  This proves membership only, not
$\exists\R$-hardness.
\end{proof}

\section{NP-hardness}
\label{sec:hardness}

We reduce from \textsc{Partition}.  The source matrix family follows
Blondel and Tsitsiklis \cite{BlondelTsitsiklis1997}; all additional lifting,
scaling, and mass-action realization steps are stated below.

Given positive integers $a_1,\ldots,a_\ell$, choose
$k=\lceil\sqrt{\ell+1}\rceil$, put $m=k^2$, and pad to $a\in\mathbb Z^m$.
Let
\[
 \gamma=a^Ta,
 \quad Q=I_m+aa^T,
 \quad \beta=1-\frac{1}{2m(1+\gamma)},
\]
and
\begin{equation}
 B(x,y)=\begin{pmatrix}-kQ&y\\x^T&k\beta\end{pmatrix},
 \qquad x,y\in(-1,1)^m.
 \label{eq:openCube}
\end{equation}

\begin{lemma}[Open-cube partition family]
\label{lem:openCube}
The family \eqref{eq:openCube} contains a Hurwitz matrix if and only if the
\textsc{Partition} instance is YES.
\end{lemma}

\begin{proof}
If $t\in\{\pm1\}^m$ and $a^Tt=0$, set
$r=(1+\beta)/2$, $x=rt$, and $y=-rt$.  The subspace spanned by $a$ has
eigenvalue $-k(1+\gamma)$, the complement orthogonal to $a,t$ has eigenvalue
$-k$, and the remaining block is
\[
 \begin{pmatrix}-k&-rk\\rk&k\beta\end{pmatrix}.
\]
Its trace is negative and determinant
$k^2(r^2-\beta)=k^2(1-\beta)^2/4>0$.

Conversely, interpolate from $(x,y)=(0,0)$, where there is exactly one positive
eigenvalue, to a Hurwitz endpoint.  The signed determinant changes sign, so
$B(\theta x,\theta y)$ is singular for some $0<\theta<1$.  With $z=-y$, the
Schur complement gives
\[
 \theta^2x^TQ^{-1}z=m\beta,
\]
hence $x^TQ^{-1}z>m\beta$.  Positive definiteness gives
\[
 \max_{x,z\in[-1,1]^m}x^TQ^{-1}z
 =\max_{t\in\{\pm1\}^m}t^TQ^{-1}t.
\]
Using $Q^{-1}=I-aa^T/(1+\gamma)$, this maximum equals
\[
 m-\frac{\min_t(a^Tt)^2}{1+\gamma}.
\]
Since $m\beta=m-1/[2(1+\gamma)]$, the integer $a^Tt$ must vanish.
\end{proof}

Write the varying entries as a rank-one affine family
$B(q)=B_0+UX(q)$, where row $r$ of $X(q)$ is $q_rw_r^T$, and define
\begin{equation}
 L(q)=\begin{pmatrix}
 B_0&U\\
 X(q)(B_0+I)&X(q)U-I
 \end{pmatrix}.
 \label{eq:lift}
\end{equation}

\begin{lemma}[Row splitting and right-scaling elimination]
\label{lem:scaling}
The lower-unitriangular matrix
$P(q)=\left(\begin{smallmatrix}I&0\\-X(q)&I\end{smallmatrix}\right)$ satisfies
\[
 P(q)L(q)P(q)^{-1}
 =\begin{pmatrix}B(q)&U\\0&-I\end{pmatrix}.
\]
Moreover,
\[
 \exists q\in(-1,1)^{2m},\ \exists D\succ0:
 L(q)D\text{ Hurwitz}
\]
if and only if the partition instance is YES.
\end{lemma}

\begin{proof}
The similarity follows by direct block multiplication.  A YES witness uses
$D=I$.  Conversely, write $D=\diag(D_0,D_1)$.  At $q=0$ the matrix is block
upper triangular.  Its $-kQD_{0,t}$ block is similar to a negative definite
matrix, the scalar $k\beta d_*$ is positive, and $-D_1$ is negative definite.
Thus $L(0)D$ has exactly one positive eigenvalue.  A Hurwitz endpoint has the
opposite signed determinant, so some $L(\theta q)D$ with $0<\theta<1$ is
singular.  Since $D$ is nonsingular and
$\det L(q)=(-1)^{2m}\det B(q)$, the original matrix $B(\theta q)$ is singular.
The strict cube argument in \cref{lem:openCube} then forces a partition.
\end{proof}

The lift has $m+1$ fixed rows and $2m$ rows of the form $a_i+q_ib_i$.

\begin{lemma}[Exact mass-action row realization]
\label{lem:rowrealization}
For any rational independent-row family $L(q)$ with fixed and affine
one-parameter rows, a classical mass-action network on the same number of
species has complete positive steady-flux image
\[
 \{\Gamma\diag(v)Y^T:v>0,\Gamma v=0\}
 =\{R L(q):R\succ0,\ q\in(-1,1)^d\}.
\]
It has full stoichiometric rank and uses two reactions per fixed row and three
per variable row.
\end{lemma}

\begin{proof}
For a variable row choose an integer $c$ clearing denominators of $a\pm b$ and
a sufficiently large integer base complex $y^0$ with the row species present.
Let $y^\pm=y^0+c(a\pm b)$ and use reactions
\[
 y^+\to y^++e_i,
 \qquad y^-\to y^-+e_i,
 \qquad y^0\to y^0-e_i.
\]
Positive steady fluxes have the exact form $u,w,u+w$.  Their row contribution
is
\[
 c\bigl((u+w)a+(u-w)b\bigr)=\rho(a+qb),
\quad
\rho=c(u+w)>0,
\quad
q=\frac{u-w}{u+w}\in(-1,1).
\]
Every $\rho,q$ is recovered by
$u=\rho(1+q)/(2c)$, $w=\rho(1-q)/(2c)$.  A fixed row uses the analogous
$+e_i,-e_i$ pair.  Because each row uses only its own stoichiometric direction,
$\Gamma v=0$ contains exactly these rowwise equations and no cross-row
constraints.  Each direction $e_i$ occurs, so the rank is full.
\end{proof}

\begin{theorem}[NP-hardness]
\label{thm:nphard}
Weak all-mobile stationary-Turing admissibility is NP-hard under
polynomial-time many-one reductions from \textsc{Partition}.  An instance of
length $\ell$ produces a full-rank network with $3m+1$ species and $8m+2$
reactions, where $m=k^2>\ell$, and every stoichiometric column is $\pm e_i$.
All rates and equilibrium coordinates in a YES witness are strictly positive.
The reaction molecularity is not uniformly bounded.
\end{theorem}

\begin{proof}
Apply \cref{lem:rowrealization} to $L(q)$.  Every positive-equilibrium Jacobian
is
\[
 J=R L(q)H,
\]
which is similar to $L(q)HR$.  By \cref{lem:scaling}, such a Jacobian is
Hurwitz exactly for YES partition instances.  The coordinate inherited from
the bottom-right entry of $B$ satisfies $L_{jj}=k\beta>0$, so every realized
Jacobian has $J_{jj}>0$ and hence a negative signed singleton minor.  A Hurwitz
realization is therefore stationary-destabilizable by \cref{thm:fixedJ}; the
reverse implication includes homogeneous stability by definition.

There are $2(m+1)+3(2m)=8m+2$ reactions.  All rational entries, denominator
clearing factors, and nonnegative base complexes have polynomial bit length.
A YES witness takes the rational open-cube point, $R=H=I$, $x^*=\mathbf1$, and
rates equal to the positive row fluxes.
\end{proof}

This theorem does not imply that no compact criterion can exist.  It implies
that polynomial-time recognition is impossible unless $\mathrm P=\mathrm{NP}$;
a compact condition could itself be computationally hard to evaluate.

\section{Polynomial-time algorithms for fixed species count}
\label{sec:fixedn}

Define
\[
 K=\{v\in\R_{\geq0}^m:\Gamma v=0\},
 \qquad
 \Phi(v)=\Gamma\diag(v)Y^T,
 \qquad
 P=\Phi(K)\subseteq\R^{n\times n}.
\]

\begin{lemma}[Positive fluxes project to the relative interior]
\label{lem:ri}
If $K$ contains a strictly positive vector, then
\[
 \{\Phi(v):v>0,\Gamma v=0\}=\ri P.
\]
\end{lemma}

\begin{proof}
The strict point ensures that no coordinate inequality is identically tight,
so $\ri K=K\cap\R_{>0}^m$.  Linear maps preserve relative interiors of
nonempty convex sets \cite{Rockafellar1970}.
\end{proof}

\begin{theorem}[Fixed-species polynomial time]
\label{thm:fixedn}
For every fixed $n$, weak all-mobile stationary-Turing admissibility of
rationally encoded classical mass-action networks is decidable in polynomial
bit complexity in the reaction count and input length.  The same holds when a
designated mobile set is included in the input.
\end{theorem}

\begin{proof}
Strict positive-flux feasibility is a rational linear program: by homogeneity,
it is equivalent to $\Gamma v=0$ and $v_r\geq1$.  Every extreme ray of $K$ is
a support-minimal nonnegative dependency among columns of $\Gamma$, hence has
support at most $\rank\Gamma+1\leq n+1$.  For fixed $n$, enumerate all such
column subsets, compute their exact one-dimensional positive kernels, and map
the resulting circuit rays through $\Phi$.  There are $O(m^{n+1})$ tests and
the generators have polynomial bit length by determinant bounds.

These images generate $P$ in fixed ambient dimension $n^2$.  Exact
fixed-dimensional polyhedral algorithms compute its span, lineality, and a
relative-facet description in polynomial time \cite{Schrijver1986}.  Thus
$A\in\ri P$ is described by linear equations and strict rational inequalities,
including lower-dimensional or nonpointed images.

By \cref{lem:ri}, the remaining formula uses only the $n^2$ entries of $A$ and
$n$ positive entries of $h$: relative-facet conditions, Routh--Hurwitz
inequalities for $A\diag(h)$ on $S$, and one negative signed principal minor.
For fixed $n$, the variable count and degrees are fixed and the number and bit
length of coefficients are polynomial.  Fixed-variable algorithms for the
first-order theory of the reals therefore run in polynomial bit complexity
\cite{Renegar1992,BasuPollackRoy2006}.  A designated mobile set adds one
spectral variable and one of constantly many block choices.
\end{proof}

The conclusion is a polynomial-time exact decision algorithm, not a finite
structural classification.  In particular it applies to arbitrary reaction
lists for $n=1,2,3$.

\section{A biochemical two-species benchmark}
\label{sec:example}

Consider the Schnakenberg reactions
\[
 0\xrightarrow{1/10}X,
 \qquad X\xrightarrow{1}0,
 \qquad 0\xrightarrow{9/10}Y,
 \qquad 2X+Y\xrightarrow{1}3X.
\]
The mass-action equations are
\[
 \dot x=\frac1{10}-x+x^2y,
 \qquad
 \dot y=\frac9{10}-x^2y.
\]
At $(x^*,y^*)=(1,9/10)$ the positive flux is
$(1/10,1,9/10,9/10)$ and the Jacobian is
\[
 J=\begin{pmatrix}4/5&1\\-9/5&-1\end{pmatrix}.
\]
Its trace is $-1/5$ and determinant is $1$, so it is Hurwitz.  The positive
entry $J_{11}$ is a negative signed singleton minor.  Indeed, with
\[
 D=\diag(1/100,10),
 \qquad
 \det(D-J)=-\frac{689}{100}<0,
\]
$J-D$ has a positive real eigenvalue.  This example illustrates the exact
network formula in the classical activator--inhibitor setting; it does not
assert a nonlinear patterned branch.

\section{Data level, designated mobility, and certificates}
\label{sec:extensions}

The exact formula uses the indexed pair $(\Gamma,Y)$.  Among the six levels
considered---unsigned interaction graph, signed interaction graph, a Coates
graph at one equilibrium, $\Gamma$ alone, $(\Gamma,Y)$, and the fully labeled
indexed complex graph---Level 5 is the weakest sufficient parameter-independent
input.  Exact separating network pairs show that Levels 1 and 2 can have the
same signs but opposite answers and that the same ordered $\Gamma$ with
different sources can have opposite answers.  Level 3 depends on a selected
parameter realization.  Levels 5 and 6 are equivalent when reaction indices
and integer labels are retained, since $y'_r=y_r+\gamma_r$.  This statement is
only relative to these six representations.

If only species in $M$ diffuse, the determinant-at-zero principal-minor rule is
false.  The correct fixed-J statement uses a positive spectral parameter.

\begin{theorem}[Designated mobile species]
\label{thm:mobile}
Assume $\det(\lambda I-J)>0$ for every $\lambda>0$.  There is a diagonal
matrix positive on $M$ and zero on $M^c$ such that $J-D$ has a positive real
eigenvalue if and only if there are $\lambda>0$ and
$I\supseteq M^c$ with
\[
 \det(\lambda I_{|I|}-J[I,I])<0.
\]
\end{theorem}

\begin{proof}
Expanding only in mobile diagonal variables gives
\[
 \det(\lambda I-J+D_M)
 =\sum_{I\supseteq M^c}
 \det(\lambda I_{|I|}-J[I,I])
 \prod_{j\in M\setminus I}d_j.
\]
A zero with positive background coefficient requires a negative coefficient.
Conversely, the same multiscale assignment used in \cref{thm:fixedJ} uniquely
dominates a selected negative coefficient; the determinant is then negative
at the chosen $\lambda$ and positive for large spectral parameter, yielding a
larger positive root.  For mass-action homogeneous stability,
\eqref{eq:positive-axis} supplies the hypothesis.
\end{proof}

An exact counterexample to the zero-parameter substitute is
\[
 J=\begin{pmatrix}-10&20&-50\\1&1&0\\1&0&2\end{pmatrix},
 \qquad M=\{1\}.
\]
The matrix is Hurwitz and the relevant zero-parameter signed block test is
nonnegative, but
$J-\diag(257/2,0,0)$ has eigenvalue $3/2$; at $\lambda=3/2$ the immobile block
determinant is $-1/4$.

\begin{theorem}[Finite exact certificates]
\label{thm:certificates}
Every YES instance has a finite exact real-algebraic sample certificate.
Every NO instance has a finite exact Real-Nullstellensatz certificate.  Both
are independently checkable.  No polynomial certificate-size bound or
practical general certificate generator is claimed.
\end{theorem}

\begin{proof}[Proof sketch]
Replace $g>0$ by $gz^2-1=0$ and $g<0$ by $(-g)z^2-1=0$.  A nonempty rational
semialgebraic set has a point with real-algebraic coordinates, representable in
one isolated number field.  If a branch has no real zero, the Real
Nullstellensatz gives
\[
 -1=\sum_j s_j^2+\sum_kq_kF_k.
\]
Exact coefficient comparison verifies the identity.  There are finitely many
principal branches.  Detailed formats and generation/size distinctions are in
the supplement \cite{BochnakCosteRoy1998,BasuPollackRoy2006}.
\end{proof}

\section{Discussion}

The exact characterization exposes the two coupled parts of the network-wide
problem: positive steady-flux geometry and positive right-diagonal homogeneous
stabilization.  The NP-hardness theorem shows that their interaction already
encodes a classical combinatorial decision problem even when every
stoichiometric column is a unit coordinate change.  For fixed species count,
the flux variables can nevertheless be projected away into fixed matrix
dimension, producing polynomial-time exact algorithms.

Several directions remain open.  Most important biologically is whether the
hardness survives bimolecular or uniformly bounded-molecularity networks.  The
present denominator-clearing construction does not.  Strong NP-hardness and
$\exists\R$-hardness are also unknown.  At the dynamical level, weak
stationary admissibility has not been shown equivalent to a primary simple
transverse stationary crossing, and no arbitrary-dimensional wave criterion
or nonlinear patterned-state theorem is supplied.

The accompanying software contains an independent constructor, verifier, and
adversarial tests.  It does not bundle a complete general-purpose
quantifier-elimination engine: a raw network without an attached exact
certificate may remain unresolved computationally even though the abstract
decision procedure terminates.

\paragraph{Computational assistance and reproducibility.}
Symbolic computation, numerical falsification, and AI-assisted drafting were
used in developing and auditing the work.  The mathematical claims in this
version are supported by explicit proofs and exact independently replayable
artifacts; numerical searches are not used as proofs.

\begingroup
\small
\bibliographystyle{plain}
\bibliography{references}
\endgroup
\end{document}
